\begin{helper}[A scale-invariant inequality on a ball passes from smooth compactly supported functions to the $W^{1,2}$ class]
Let $n\ge2$ and let $q,a,C_1\in\R$ with $q>1$, $a\ge0$ and $C_1>0$.  Assume the
inequality is known in the smooth case:
\begin{enumerate}
\item[(H)] for every $x_0\in\R^n$, every $\rho>0$ and every
  $\Phi\in C^\infty(\R^n)$ with compact support and
  $\operatorname{supp}\Phi\subseteq B_\rho(x_0)$,
  \[
  \Bigl(\int_{B_\rho(x_0)}|\Phi|^{q}\,dx\Bigr)^{1/q}
  \ \le\ C_1\rho^{a}
  \Bigl(\int_{B_\rho(x_0)}|\nabla\Phi|^2\,dx\Bigr)^{1/2},
  \]
  where $\nabla\Phi(x)$ is the vector of the partial derivatives
  $D_i\Phi(x)=d\Phi_x(e_i)$ and $|\cdot|$ its Euclidean norm.
\end{enumerate}
Then the same inequality holds for the $W^{1,2}$ class: for every
$x_0\in\R^n$, every $\rho>0$ and every pair $(\varphi,g_\varphi)$ with
$\varphi\in W^{1,2}(B_\rho(x_0))$ of weak gradient $g_\varphi$
(\noderef{MemW1pOn}), $\varphi$ of compact support and
$\operatorname{supp}\varphi\subseteq B_\rho(x_0)$,
\[
\Bigl(\int_{B_\rho(x_0)}|\varphi|^{q}\,dx\Bigr)^{1/q}
\ \le\ C_1\rho^{a}
\Bigl(\int_{B_\rho(x_0)}|g_\varphi|^2\,dx\Bigr)^{1/2}.
\]
\end{helper}
\begin{proof}
Write $B_\rho=B_\rho(x_0)$.  Since $q>1$ we have $q>0$, so $t\mapsto t^{1/q}$
is nondecreasing on $[0,\infty)$ and $t\mapsto t^{q}$ is nonnegative; these are
used at the end.

\emph{Step 1: the smooth approximating sequence.}  Apply
\noderef{WeakGradientSmoothApprox} with this $n$, the exponent $2$, the open set
$B_\rho$ and the pair $(\varphi,g_\varphi)$; its hypotheses are $1\le n$ (from
$n\ge2$), $1\le2$, openness of $B_\rho$, the $W^{1,2}$ membership, compactness
of the support of $\varphi$ and $\operatorname{supp}\varphi\subseteq B_\rho$.
It produces a compact $K$ with
$\operatorname{supp}\varphi\subseteq K\subseteq B_\rho$ and functions
$\Phi_m\in C^\infty(\R^n)$, each with compact support contained in $K$, such
that
\[
\int_K|\Phi_m-\varphi|^2\,dx\to0
\qquad\text{and}\qquad
\int_K\bigl|D_i\Phi_m-(g_\varphi)_i\bigr|^2\,dx\to0
\quad(m\to\infty)
\tag{1}
\]
for every index $i$.  In particular $K$ has finite measure, being compact.

\emph{Step 2: the smooth case at each $\Phi_m$.}  Each $\Phi_m$ is $C^\infty$
with compact support contained in $K\subseteq B_\rho$, hence
$\operatorname{supp}\Phi_m\subseteq B_\rho$, so (H) applies to it at this $x_0$
and this $\rho$:
\[
\Bigl(\int_{B_\rho}|\Phi_m|^{q}\Bigr)^{1/q}
\le C_1\rho^{a}\Bigl(\int_{B_\rho}|\nabla\Phi_m|^2\Bigr)^{1/2}.
\tag{2}
\]
We pass to the limit in (2).

\emph{Step 3: which integrands vanish off $K$.}  The functions $\varphi$ and
each $\Phi_m$ vanish off $K$, their closed supports lying in $K$; and
$\nabla\Phi_m$ vanishes off $K$ as well, because the differential of a function
vanishes at every point outside its closed support, so
$\nabla\Phi_m(x)=0$ for $x\notin\operatorname{supp}\Phi_m$ and a fortiori for
$x\notin K$.  Since $K\subseteq B_\rho$, for each of these three integrands the
integral over $B_\rho$ equals the integral over $K$:
\[
\int_{B_\rho}|\Phi_m|^{q}=\int_K|\Phi_m|^{q},
\qquad
\int_{B_\rho}|\varphi|^{q}=\int_K|\varphi|^{q},
\qquad
\int_{B_\rho}|\nabla\Phi_m|^{2}=\int_K|\nabla\Phi_m|^{2}.
\tag{2a}
\]
\emph{No such claim is made about $g_\varphi$}, which need not vanish off $K$;
for it we use only that the integrand $|g_\varphi|^2$ is nonnegative and
$K\subseteq B_\rho$, so that by monotonicity of the integral (Preamble)
\[
\int_K|g_\varphi|^{2}\ \le\ \int_{B_\rho}|g_\varphi|^{2}=:G^{2},
\qquad G\ge0 ,
\tag{2b}
\]
and $G<\infty$ because $|g_\varphi|\in L^2(B_\rho)$ by \noderef{MemW1pOn}.

\emph{Step 4: the right-hand side.}  Put
\[
\epsilon_m=\sum_{i=1}^n
\Bigl(\int_K\bigl|D_i\Phi_m-(g_\varphi)_i\bigr|^{2}\Bigr)^{1/2}\ \ge\ 0 ,
\]
so $\epsilon_m\to0$ by (1).  By the $L^2$ triangle inequality (Minkowski,
Preamble), then $|v|\le\sum_i|v_i|$ for $v\in\R^n$ applied pointwise to
$v=\nabla\Phi_m-g_\varphi$ together with Minkowski again, and finally (2b),
\[
\Bigl(\int_K|\nabla\Phi_m|^{2}\Bigr)^{1/2}
\le\Bigl(\int_K|g_\varphi|^{2}\Bigr)^{1/2}
+\Bigl(\int_K|\nabla\Phi_m-g_\varphi|^{2}\Bigr)^{1/2}
\le G+\epsilon_m .
\]
With (2a) and $C_1\rho^{a}\ge0$, inequality (2) therefore reads
\[
\Bigl(\int_K|\Phi_m|^{q}\Bigr)^{1/q}
\ \le\ C_1\rho^{a}\bigl(G+\epsilon_m\bigr)
\qquad(m\in\mathbb{N}).
\tag{4}
\]

\emph{Step 5: the left-hand side.}  By the first limit in (1) and the Preamble
item on extracting an almost everywhere convergent subsequence from $L^2$
convergence --- applicable because $K$ has finite measure, being compact, and
the exponent $2$ is finite --- there is a subsequence $\Phi_{m_l}$ with
$\Phi_{m_l}\to\varphi$ almost everywhere on $K$.  The functions
$|\Phi_{m_l}|^{q}$ are nonnegative and measurable and converge almost
everywhere on $K$ to $|\varphi|^{q}$, so Fatou's lemma (Preamble) gives
\[
\int_K|\varphi|^{q}\ \le\ \liminf_{l\to\infty}\int_K|\Phi_{m_l}|^{q} .
\tag{5}
\]

\emph{Step 6: conclusion.}  Raising (4) to the power $q$ --- legitimate because
both sides are nonnegative and $t\mapsto t^{q}$ is nondecreasing on
$[0,\infty)$ --- gives
$\int_K|\Phi_{m}|^{q}\le\bigl(C_1\rho^{a}(G+\epsilon_m)\bigr)^{q}$ for every
$m$.  Along the subsequence the right-hand side converges to
$\bigl(C_1\rho^{a}G\bigr)^{q}$, since $\epsilon_{m_l}\to0$, $G<\infty$ and
$t\mapsto(C_1\rho^{a}t)^q$ is continuous on $[0,\infty)$.  With (5),
\[
\int_K|\varphi|^{q}\ \le\ \bigl(C_1\rho^{a}G\bigr)^{q}<\infty .
\]
By (2a) the left-hand side is $\int_{B_\rho}|\varphi|^{q}$, so that integral is
\emph{finite} --- the integral on the left of the assertion is a genuine finite
integral and not the value a divergent one would be assigned.  Finally
$t\mapsto t^{1/q}$ is nondecreasing on $[0,\infty)$ and both sides are
nonnegative, so taking $q$-th roots and recalling
$G=\bigl(\int_{B_\rho}|g_\varphi|^2\bigr)^{1/2}$,
\[
\Bigl(\int_{B_\rho}|\varphi|^{q}\Bigr)^{1/q}
\le C_1\rho^{a}\Bigl(\int_{B_\rho}|g_\varphi|^2\Bigr)^{1/2},
\]
which is the assertion.
\end{proof}
